\PageHeader
  {Appendix: aceleași idei în C++}
  {Până acum am lucrat în pseudocod. Acum vedem cum arată aceleași concepte într-un limbaj real.}
  {Limbaje 1 din 3}

\begin{languagebox}[Ce să observi]
În C++ apar tipurile, acoladele și punctul și virgula. Ideile rămân aceleași: variabile, modificări de valori, condiții, repetiții.
\end{languagebox}

\begin{examplebox}[Variabile și \texttt{if}]
\begin{lstlisting}[style=languagequest,language=C++]
int score = 0;
int lives = 3;

score = score + 5;
lives = lives - 1;

if (score >= 10) {
    cout << "Ai trecut mai departe.";
}
\end{lstlisting}
\end{examplebox}

\begin{examplebox}[\texttt{while} și \texttt{for}]
\begin{lstlisting}[style=languagequest,language=C++]
int i = 1;
while (i <= 5) {
    cout << i << " ";
    i = i + 1;
}

for (int pas = 1; pas <= 3; pas++) {
    cout << pas << " ";
}
\end{lstlisting}
\end{examplebox}

\begin{languagebox}[Ne amintim]
C++ cere mai multă grijă la sintaxă, dar logica programului este aceeași pe care ai exersat-o deja în pseudocod.
\end{languagebox}

\clearpage

\PageHeader
  {Appendix: aceleași idei în Python}
  {Python scrie mai puțin și lasă indentarea să arate structura programului.}
  {Limbaje 2 din 3}

\begin{languagebox}[Ce să observi]
În Python nu mai declarăm tipul în fața variabilei, iar blocurile sunt marcate prin indentare, nu prin acolade.
\end{languagebox}

\begin{examplebox}[Variabile și \texttt{if}]
\begin{lstlisting}[style=languagequest,language=Python]
score = 0
lives = 3

score = score + 5
lives = lives - 1

if score >= 10:
    print("Ai trecut mai departe.")
\end{lstlisting}
\end{examplebox}

\begin{examplebox}[\texttt{while} și \texttt{for}]
\begin{lstlisting}[style=languagequest,language=Python]
i = 1
while i <= 5:
    print(i, end=" ")
    i = i + 1

for pas in range(1, 4):
    print(pas, end=" ")
\end{lstlisting}
\end{examplebox}

\begin{languagebox}[Ne amintim]
Python este foarte bun pentru a vedea logic fluxul programului, tocmai pentru că scrie concis.
\end{languagebox}

\clearpage

\PageHeader
  {Appendix: aceleași idei în Java}
  {Java seamănă la structură cu C++, dar este mai strict și mai explicit în multe locuri.}
  {Limbaje 3 din 3}

\begin{languagebox}[Ce să observi]
În Java vei vedea tipuri, acolade și instrucțiuni terminate cu punct și virgulă, la fel ca în C++. În schimb, afișarea și organizarea codului sunt puțin diferite.
\end{languagebox}

\begin{examplebox}[Variabile și \texttt{if}]
\begin{lstlisting}[style=languagequest,language=Java]
int score = 0;
int lives = 3;

score = score + 5;
lives = lives - 1;

if (score >= 10) {
    System.out.println("Ai trecut mai departe.");
}
\end{lstlisting}
\end{examplebox}

\begin{examplebox}[\texttt{while} și \texttt{for}]
\begin{lstlisting}[style=languagequest,language=Java]
int i = 1;
while (i <= 5) {
    System.out.print(i + " ");
    i = i + 1;
}

for (int pas = 1; pas <= 3; pas++) {
    System.out.print(pas + " ");
}
\end{lstlisting}
\end{examplebox}

\clearpage
